\section{解三角形}
解三角形, 就是给定若干三角形的元素, 确定出三角形其他元素的过程. 一个自然的问题是, 究竟
什么样的条件才能唯一确定一个三角形, 或者退一步, 允许一两个多解的存在？
在初中阶段就学习了全等三角形的判定条件：$SSS,SAS,ASA,AAS,HL$. 也就是说
当给出这些条件时, 应当能唯一一个三角形. 而$SSA$这样的条件, 虽然不能确定唯一一个三角形, 
但是解也不会超过两个. 但这些仅是定性地判断存在性, 如何具体计算它们的数值呢？

例如, 理论上给出两边及其夹角, 其余的一角两边应当随之确定. 可是要如何计算呢？
这些问题就由正弦定理和余弦定理来解决. 

\subsection{正弦定理}
\begin{example}
    给定一个$\rt \triangle ABC$, \(\angle C =\dfrac{\pi}{2}\), \(\angle A=\dfrac{\pi}{6}\), 
    \(AB=2\). 请尝试确定三角形其余一角和两边, 以及外接圆的半径. 
\end{example}
\begin{jie}
    
\end{jie}
\begin{theorem}[正弦定理]
  \begin{equation}
      \frac{a}{\sin A}=\frac{b}{\sin B }=\frac{c}{\sin C}=2R
  \end{equation}
\end{theorem}
\clearpage


\subsection{余弦定理}
\begin{theorem}[余弦定理]
\begin{equation}
    a^2=b^2+c^2-2bc\cos A
\end{equation}
\begin{equation}
    b^2=a^2+c^2-2ac\cos B
\end{equation}
\begin{equation}
    c^2=a^2+b^2-2ab\cos C
\end{equation}
\end{theorem}
\clearpage

\subsection{正余弦定理的综合应用}
\begin{example}
    在\(\tri ABC\)中, \(AB=1\), \(AC=2\), \(\angle BAC=\theta\). 再取一点\(D\), 使得
    \(\tri BCD\)为等边三角形. 求线段\(AD\)长的最大值. 
\end{example}
\begin{analyse}
    要求\(AD\)的长度, 总是避不开\(\angle ACB\)或者\(\angle ABC\), 无法直接用\(\theta\). 但是
    当两边及其夹角确定的时候, 整个三角形应该是唯一确定的, 故可以用已有的两边及其夹角表示出其他量. 
\end{analyse}
\begin{jie}
    通过绘图可知\(\theta\in(0,\pi)\), 最后\(AD\)的长度应当是\(\theta\)的函数. 
    不妨设\(\angle BCA=\varphi\)

    在\(\tri ABC\)中, 由正弦定理得\[\frac{BC}{\sin\theta}=\frac{AB}{\sin\varphi}\]
    即\[\sin\varphi=\frac{\sin\theta}{BC}\]

    在\(\tri ABC\)中, 由射影定理得\[AC=AB\cos\theta+BC\cos\varphi\]
    即\[\cos\varphi=\frac{2-\cos\theta}{BC}\]

    在\(\tri ABC\)中, 由余弦定理得
    \begin{align*}
         BC^2&=AC^2+AB^2-2AB\cdot AC\cos\theta\\
            &=5-4\cos\theta
    \end{align*}
   
    在\(\tri ACD\)中, $CD=BC$, 再由余弦定理得
    \begin{align*}
        AD^2&=AC^2+CD^2-2AC\cdot CD\cos(\varphi+\frac{\pi}{3})\\
            &=4+BC^2-4BC(\cos\varphi\cos\frac{\pi}{3}-\sin\varphi\sin\frac{\pi}{3})\\
            &=4+5-4\cos\theta-4(\frac{2-\cos\theta}{2}-\frac{\sqrt{3}\sin\theta}{2})\\
            &=5+2\sqrt{3}\sin\theta-2\cos\theta\\
            &=5+\sqrt{(2\sqrt{3})^2+2^2}(\frac{\sqrt{3}}{2}\sin\theta-\frac{1}{2}\cos\theta)\\
            &=5+4\sin(\theta-\frac{\pi}{6})\\
            &\leq9
    \end{align*}
    等号当\(\theta=\dfrac{2}{3}\pi\)时取得

\end{jie}


\begin{example}
    在 $ \triangle A B C$  中, 角 $ A, B, C$  所对的边分别为$  a, b, c $ , 
    且外接圆半径为 $ R=5 $ , 求$ \dfrac{a b c}{a^{2}+b^{2}+2 c^{2}} $ 的最大值为   
\end{example}

\begin{jie}
    \[
\begin{aligned}
    \frac{abc}{a^2 + b^2 + 2c^2} &= \frac{ab \cdot 2 \times 5 \sin C}{a^2 + b^2 + 2(a^2 + b^2 - 2ab \cos C)}\\
                                & = \frac{10ab \sin C}{3(a^2 + b^2) - 4ab \cos C}\\
                               &  \leq \frac{10ab \sin C}{6ab - 4ab \cos C} = \frac{5 \sin C}{3 - 2 \cos C} \\
                               &\leq \frac{5 \sin C}{\sqrt{5} \sin C} = \sqrt{5}.
\end{aligned}
\]
\qed
\end{jie}
\begin{theorem}[Weisenböck不等式]
    在\(\tri ABC\) 中, 角 $ A, B, C$  所对的边分别为$  a, b, c $ , 三角形面积为\(S\)对任意实数 $ x, y, z $ , 满足 $ x+y+z>0,  x y+y z+z x>0$  , 则有

$$x a^{2}+y b^{2}+z c^{2} \geqslant 4 \sqrt{x y+y z+z x} S$$


当且仅当 $ \cot A: \cot B: \cot C=x: y: z$  时取得等号．
\end{theorem}
\begin{proof}
    \[\begin{aligned}
        x a^{2}+y b^{2}+z c^{2} & =x a^{2}+y b^{2}+z\left(a^{2}+b^{2}-2 a b \cos C\right) \\
        & =(x+z) a^{2}+(y+z) b^{2}-2 z a b \cos C \\
        & \geqslant 2 \sqrt{(x+z)(y+z)} a b-2 z a b \cos C \\
        & =2 a b(\sqrt{(x+z)(y+z)}-z \cos C) \\
        & \geqslant 2 a b\left(\sqrt{(x+z)(y+z)-z^{2}} \sin C\right) \\
        & =4 S_{\triangle} \sqrt{x y+y z+z x}
        \end{aligned}\]

        当且仅当 $ \frac{a^{2}}{x+y}=\frac{b^{2}}{y+z}=\frac{c^{2}}{z+x} $ 亦即 $ x: y: z=\cot A: \cot B: \cot C $ 时取得等号．
    \[m \sin \theta+n \cos \theta \leqslant \sqrt{m^{2}+n^{2}} \Longleftrightarrow \sqrt{m^{2}+n^{2}}-n \cos \theta \geqslant m \sin \theta\]
        \qed
\end{proof}
\clearpage
\review
\clearpage

\hw
\begin{homework}
在$\triangle ABC$中, $a=2$, $A=\frac{\pi}{6}$, 则“$B=\frac{\pi}{3}$”是“$b=2\sqrt{3}$”的
\begin{table}[H]
  \centering
\begin{tabular*}{\hsize}{@{}@{\extracolsep{\fill}}llll@{}}
  \text{A.} 充分而不必要条件   &   \text{B.} 必要而不充分条件   &      \text{C.} 充分必要条件   &  \text{D.} 既不充分也不必要条件
\end{tabular*}
\end{table}
\end{homework}

\begin{homework}[2023年北京卷7]
在$\triangle ABC$中, $(a+c)(\sin A-\sin C)=b(\sin A-\sin B)$, 则$\angle C=$
\begin{table}[H]
  \centering
\begin{tabular*}{\hsize}{@{}@{\extracolsep{\fill}}llll@{}}
  \text{A.} $\frac{\pi}{6}$   &   \text{B.} $\frac{\pi}{3}$   &      \text{C.} $\frac{2\pi}{3}$   &  \text{D.} $\frac{5\pi}{6}$
\end{tabular*}
\end{table}
\end{homework}

\begin{homework}[2011北京高考（理）9]
在$\triangle ABC$中, 若$b=5$, $\angle B=\frac{\pi}{4}$, $\tan A=2$, 则$\sin A=$\rule{1cm}{0.15mm}; $a=$\rule{1cm}{0.15mm}.
\end{homework}

\begin{homework}[2015北京高考12]
在$\triangle ABC$中, $a=4$, $b=5$, $c=6$, 则$\frac{\sin 2A}{\sin C}=$\rule{1cm}{0.15mm}.
\end{homework}

\begin{homework}[2018北京高考（文）14]
若$\triangle ABC$的面积为$\frac{\sqrt{3}}{4}(a^{2}+c^{2}-b^{2})$, 且$\angle C$为钝角, 则$\angle B=$\rule{1cm}{0.15mm}; $\frac{c}{a}$的取值范围是\rule{1cm}{0.15mm}.
\end{homework}

\begin{homework}[2022北京高考16]
在$\triangle ABC$中, $\sin 2C = \sqrt{3} \sin C$.
\begin{enumerate}
    \item 求$\angle C$;
    \item 若$b=6$, 且$\triangle ABC$的面积为$6\sqrt{3}$, 求$\triangle ABC$的周长.
\end{enumerate}


\end{homework}

\begin{homework}[2014北京高考（理）15]
如图, 在$\triangle ABC$中, $\angle B=\frac{\pi}{3}$, $AB=8$, 点$D$在$BC$边上, 且$CD=2$, $\cos \angle ADC=\frac{1}{7}$.

\begin{minipage}[H]{0.4\textwidth}
    \begin{enumerate}
        \item 求$\sin \angle BAD$;
        \item 求$BD$, $AC$的长.
    \end{enumerate}
\end{minipage}
\hfill
\begin{minipage}[H]{0.5\textwidth}
%图在test里的第一个
\end{minipage}

\end{homework}

\begin{homework}[2013北京高考理15]
在$\triangle ABC$中, $a=3$, $b=2\sqrt{6}$, $\angle B=2\angle A$.
\begin{enumerate}
    \item 求$\cos A$的值;
    \item 求$c$的值.
\end{enumerate}

\end{homework}

\begin{homework}[2019北京高考理15]
在$\triangle ABC$中, $a=3$, $b-c=2$, $\cos B=-\frac{1}{2}$.
\begin{enumerate}
    \item 求$b$,$c$的值;
    \item 求$\sin(B-C)$的值.
\end{enumerate}
\end{homework}

\begin{homework}
在$\triangle ABC$中, $2a\cos A=b\cos C+c\cos B$, 则$\angle A=$
\begin{table}[H]
  \centering
\begin{tabular*}{\hsize}{@{}@{\extracolsep{\fill}}llll@{}}
  \text{A.} $\frac{\pi}{6}$   &   \text{B.} $\frac{\pi}{3}$   &      \text{C.} $\frac{\pi}{2}$   &  \text{D.} $\frac{2\pi}{3}$
\end{tabular*}
\end{table}
\end{homework}

\begin{homework}
在$\triangle ABC$中, $a\cos B-\frac{\sqrt{3}}{2}b=c$, 则$\angle A=$
\begin{table}[H]
  \centering
\begin{tabular*}{\hsize}{@{}@{\extracolsep{\fill}}llll@{}}
  \text{A.} $\frac{\pi}{6}$   &   \text{B.} $\frac{\pi}{3}$   &      \text{C.} $\frac{2\pi}{3}$   &  \text{D.} $\frac{5\pi}{6}$
\end{tabular*}
\end{table}
\end{homework}

\begin{homework}
在$\triangle ABC$中, $\sqrt{3}a=2b\sin A$, 则$\angle B=$
\begin{table}[H]
  \centering
\begin{tabular*}{\hsize}{@{}@{\extracolsep{\fill}}llll@{}}
  \text{A.} $\frac{\pi}{6}$   &   \text{B.} $\frac{\pi}{6} \text{ 或 } \frac{5\pi}{6}$   &      \text{C.} $\frac{\pi}{3}$   &  \text{D.} $\frac{\pi}{3} \text{ 或 } \frac{2\pi}{3}$
\end{tabular*}
\end{table}
\end{homework}

\begin{homework}
在$\triangle ABC$中, $\angle B=60^{\circ}$, $b=\sqrt{7}$, $a-c=2$, 则$\triangle ABC$的面积为
\begin{table}[H]
  \centering
\begin{tabular*}{\hsize}{@{}@{\extracolsep{\fill}}llll@{}}
  \text{A.} $\frac{3\sqrt{3}}{2}$   &   \text{B.} $\frac{3\sqrt{3}}{4}$   &      \text{C.} $\frac{3}{2}$   &  \text{D.} $\frac{3}{4}$
\end{tabular*}
\end{table}
\end{homework}
\begin{homework}
在$\triangle ABC$中，若$a=\sqrt{2}$，$\angle A=\frac{\pi}{6}$，$\cos C=\frac{1}{3}$，则$c=$
\begin{table}[H]
  \centering
\begin{tabular*}{\hsize}{@{}@{\extracolsep{\fill}}llll@{}}
  \text{A.} $\frac{\sqrt{3}}{3}$   &   \text{B.} $\frac{2}{3}$   &      \text{C.} $\frac{8\sqrt{3}}{9}$   &  \text{D.} $\frac{8}{3}$
\end{tabular*}
\end{table}
\end{homework}

\begin{homework}
在$\triangle ABC$中，$\sin B=\sin 2A$，$c=2a$，则
\begin{table}[H]
  \centering
\begin{tabular*}{\hsize}{@{}@{\extracolsep{\fill}}llll@{}}
  \text{A.} $\angle B\text{为直角}$   &   \text{B.} $\angle B\text{为钝角}$   &      \text{C.} $\angle C\text{为直角}$   &  \text{D.} $\angle C\text{是钝角}$
\end{tabular*}
\end{table}

\end{homework}

\begin{homework}
在$\triangle ABC$中，$A=\frac{\pi}{4}$，$C=\frac{7\pi}{12}$，\(b=\sqrt{2}\),则$a=$
\begin{table}[H]
  \centering
\begin{tabular*}{\hsize}{@{}@{\extracolsep{\fill}}llll@{}}
  \text{A.} $\frac{\sqrt{6}}{2}$   &   \text{B.} $\frac{\sqrt{6}}{4}$   &      \text{C.} $\frac{\sqrt{3}}{2}$   &  \text{D.} $\frac{\sqrt{3}}{4}$
\end{tabular*}
\end{table}
\end{homework}

\begin{homework}
在$\triangle ABC$中，$AB=4$，$AC=5$，$\cos C=\frac{3}{4}$，则$BC$的长为
\begin{table}[H]
  \centering
\begin{tabular*}{\hsize}{@{}@{\extracolsep{\fill}}llll@{}}
  \text{A.} $3$   &   \text{B.} $4$   &      \text{C.} $5$   &  \text{D.} $6$
\end{tabular*}
\end{table}
\end{homework}

\begin{homework}
在$\triangle ABC$中，$\angle C=60^{\circ}$，$a+2b=8$，$\sin A=6\sin B$，则$c=$
\begin{table}[H]
  \centering
\begin{tabular*}{\hsize}{@{}@{\extracolsep{\fill}}llll@{}}
  \text{A.} $2\sqrt{3}$   &   \text{B.} $4$   &      \text{C.} $4\sqrt{3}$   &  \text{D.} $6$
\end{tabular*}
\end{table}
\end{homework}

\begin{homework}
在$\triangle ABC$中，若$a=2$，$b=3$，$\cos(A+B)=\frac{1}{3}$，则$c=$
\begin{table}[H]
  \centering
\begin{tabular*}{\hsize}{@{}@{\extracolsep{\fill}}llll@{}}
  \text{A.} $\sqrt{17}$   &   \text{B.} $4$   &      \text{C.} $\sqrt{15}$   &  \text{D.} $3$
\end{tabular*}
\end{table}
\end{homework}

\begin{homework}
在$\triangle ABC$中，若$c=4$，$b=a=1$，$\cos C=\frac{1}{4}$，则$\triangle ABC$的面积是
\begin{table}[H]
  \centering
\begin{tabular*}{\hsize}{@{}@{\extracolsep{\fill}}llll@{}}
  \text{A.} $1$   &   \text{B.} $\frac{3}{4}$   &      \text{C.} $\sqrt{17}$   &  \text{D.} $\frac{3\sqrt{15}}{4}$
\end{tabular*}
\end{table}
\end{homework}

\begin{homework}
在$\triangle ABC$中，$a=2\sqrt{6}$，$b=2\sqrt{6}$，$\cos A=-\frac{1}{4}$，则$\triangle ABC$的面积为
\begin{table}[H]
  \centering
\begin{tabular*}{\hsize}{@{}@{\extracolsep{\fill}}llll@{}}
  \text{A.} $\frac{3\sqrt{15}}{2}$   &   \text{B.} $4$   &      \text{C.} $\sqrt{15}$   &  \text{D.} $2\sqrt{15}$
\end{tabular*}
\end{table}
\end{homework}

\begin{homework}
在$\triangle ABC$中，$\angle A=\frac{\pi}{4}$，则“$\sin B<\frac{\sqrt{2}}{2}$”是“$\triangle ABC$是钝角三角形”的
\begin{table}[H]
  \centering
\begin{tabular*}{\hsize}{@{}@{\extracolsep{\fill}}llll@{}}
  \text{A.} 充分而不必要条件   &   \text{B.} 必要而不充分条件   &      \text{C.} 充分必要条件   &  \text{D.} 既不充分也不必要条件
\end{tabular*}
\end{table}
\end{homework}

\begin{homework}
在$\triangle ABC$中，若$2\cos A \sin B = \sin C$，则该三角形的形状一定是
\begin{table}[H]
  \centering
\begin{tabular*}{\hsize}{@{}@{\extracolsep{\fill}}llll@{}}
  \text{A.} 等腰三角形   &   \text{B.} 等边三角形   &      \text{C.} 直角三角形   &  \text{D.} 等腰直角三角形
\end{tabular*}
\end{table}
\end{homework}

\begin{homework}
在$\triangle ABC$中，“$\sin 2A = \sin 2B$”是“$\triangle ABC$为等腰三角形”的
\begin{table}[H]
  \centering
\begin{tabular*}{\hsize}{@{}@{\extracolsep{\fill}}llll@{}}
  \text{A.} 充分而不必要条件   &   \text{B.} 必要而不充分条件   &      \text{C.} 充分必要条件   &  \text{D.} 既不充分也不必要条件
\end{tabular*}
\end{table}
\end{homework}

\begin{homework}
已知$\angle A$，$\angle B$是$\triangle ABC$的内角，“$\triangle ABC$为锐角三角形”是“$\sin A > \cos B$”的
\begin{table}[H]
  \centering
\begin{tabular*}{\hsize}{@{}@{\extracolsep{\fill}}llll@{}}
  \text{A.} 充分而不必要条件   &   \text{B.} 必要而不充分条件   &      \text{C.} 充分必要条件   &  \text{D.} 既不充分也不必要条件
\end{tabular*}
\end{table}
\end{homework}

\begin{homework}
在$\triangle ABC$中，若$b=5$，$\angle B=\frac{\pi}{4}$，$\cos A=\frac{3}{5}$，则$a=$\rule{1cm}{0.15mm}.
\end{homework}

\begin{homework}
在$\triangle ABC$中，若$a=2$，$\tan A=\frac{4}{3}$，$\cos B=\frac{4}{5}$，则$b=$\rule{1cm}{0.15mm}.
\end{homework}

\begin{homework}
在$\triangle ABC$中，若$a=5$，$\angle B=\frac{\pi}{4}$，$\cos A=\frac{3}{5}$，则$a=$\rule{1cm}{0.15mm}.
\end{homework}

\begin{homework}
在$\triangle ABC$中，若$a=2$，$\tan A=\frac{4}{3}$，$\cos B=\frac{4}{5}$，则$b=$\rule{1cm}{0.15mm}.
\end{homework}

\begin{homework}
在$\triangle ABC$中，角$A$，$B$，$C$所对的边分别为$a$，$b$，$c$，且$b\sin C+\sqrt{3}c\cos B=2c$。
\begin{enumerate}
    \item 求$\angle B$；
    \item 若$a=2\sqrt{5}$，$b+c=4$，求$\triangle ABC$的面积。
\end{enumerate}

\end{homework}

\begin{homework}
在锐角$\triangle ABC$中，角$A$，$B$，$C$的对边分别为$a$，$b$，$c$，且$2b\sin A-\sqrt{3}a=0$。
\begin{enumerate}
    \item 求角$B$的大小；
    \item 求$\cos A+\cos C$的取值范围。
\end{enumerate}

\end{homework}

\begin{homework}
在$\triangle ABC$中，$a\cos C+c\cos A=\frac{2\sqrt{3}}{3}b\cos B$。
\begin{enumerate}
    \item 求$\angle B$；
    \item 若$a=12$，$D$为$BC$边的中点，且$AD=3$，求$b$的值。
\end{enumerate}

\end{homework}